%% kapa: exam=1 round=34 year=2023 period=2 form=A subject=accounting subject_ko=회계학 seq=31 number=71
%% source: Q-Net 감정평가사 기출문제 (cst00309) artlSeq=5211789 file=raw/1008/5211789/2023년도 제34회 감정평가사 1차시험 2교시.pdf
%% conversion: kordoc 4.24.51 -> tools/gen1.py (auto); review=proofread-sonnet; flags=none
\begin{question}{71}
범용기계장치를 이용하여 제품 X와 Y를 생산·판매하는 (주)감평의 당기 예산 자료는 다음과 같다.
\par{\small\centering\begin{tabular}{|c|c|c|}\hline
구 분 & 제품 X & 제품 Y \\ \hline
단위당 판매가격 & ￦1,500 & ￦1,000 \\ \hline
단위당 변동원가 & 1,200 & 800 \\ \hline
단위당 기계가동시간 & 2시간 & 1시간 \\ \hline
연간 정규시장 판매수량 & 300단위 & 400단위 \\ \hline
연간 최대기계가동시간 & \multicolumn{2}{c|}{1,000시간} \\ \hline
\end{tabular}\par}
(주)감평은 신규거래처로부터 제품 Z 200단위의 특별주문을 요청받았다. 제품 Z의 생산에는 단위당 ￦900의 변동원가가 발생하며 단위당 1.5 기계가동시간이 필요하다. 특별주문 수락 시 기존 제품의 정규시장 판매를 일부 포기해야 하는 경우, (주)감평이 제시할 수 있는 단위당 최소판매가격은? (단, 특별주문은 전량 수락하든지 기각해야 한다.)
\begin{choices}
  \item ￦900
  \item ￦1,125
  \item ￦1,150
  \item ￦1,200
  \item ￦1,350
\end{choices}
\end{question}
\answer{2}
